\documentclass[12pt, a4paper]{article}
\usepackage[UTF8]{ctex}
\usepackage{booktabs}
\usepackage{geometry}
\usepackage{array}
\usepackage{xcolor}
\usepackage{hyperref}
\usepackage{amsmath}

\geometry{margin=2.5cm}

\title{\textbf{剧透检测系统准确率测试报告}\\
\large 《球状闪电》评论剧透识别评估}
\author{CommentMask / SpoilerShield}
\date{2026 年 4 月 23 日}

\begin{document}

\maketitle
\tableofcontents
\newpage

%------------------------------------------------------------
\section{测试概述}

本报告针对书籍《球状闪电》（book\_id = 9，共 2422 段），
对剧透检测系统进行系统性准确率评估。
测试数据由 Claude claude-sonnet-4-6 批量生成，检测由本系统 RAG + LLM 流水线执行。

%------------------------------------------------------------
\section{测试数据集}

\begin{table}[h]
\centering
\begin{tabular}{lrr}
\toprule
\textbf{类别} & \textbf{预期数量} & \textbf{实际生成} \\
\midrule
剧透评论       & 300  & 292  \\
非剧透评论     & 700  & 670  \\
\midrule
合计           & 1000 & 962  \\
生成失败（跳过）& —   & 38   \\
\bottomrule
\end{tabular}
\caption{测试集构成}
\end{table}

评论长度控制在 20--50 字（中文），从书籍前半部分（1211 段）随机抽取段落生成。

%------------------------------------------------------------
\section{检测流水线}

\subsection{RAG 召回策略}

\begin{enumerate}
  \item \textbf{语义向量检索}（top-k = 5）：使用 \texttt{all-MiniLM-L6-v2} 模型对评论编码，
        与段落预计算 embedding 做余弦相似度排序，最低准入阈值 $\text{similarity} \geq 0.45$。
  \item \textbf{关键词频次检索}（top-k = 3）：统计中文字符在段落中的出现频次，
        取频次最高的 3 段作为补充召回。
  \item \textbf{合并去重}：两路结果合并，按 \texttt{para\_index} 去重。
  \item \textbf{未来段落过滤}：仅保留 \texttt{para\_index > current\_para\_index} 的段落作为剧透上下文。
  \item \textbf{上下文长度控制}：若未来段落总字数超过 1200 字，
        则对每段提取关键词密度最高的 120 字片段（Snippet）传入 LLM。
  \item \textbf{无上下文回退}：若两路召回均未通过准入阈值，
        告知 LLM「未发现相关后续内容，请仅根据常识判断」。
\end{enumerate}

\subsection{LLM 判断}

召回内容连同当前段落、评论文本，以 JSON 格式提示词发送至
\texttt{claude-sonnet-4-6}，要求返回：

\begin{itemize}
  \item \texttt{is\_spoiler}：布尔值
  \item \texttt{confidence}：0.0--1.0 置信度
  \item \texttt{reason}：判断依据
\end{itemize}

置信度阈值 $\geq 0.7$ 时标记为 \texttt{HIDDEN}（剧透）或 \texttt{PUBLIC}（公开），
否则标记为 \texttt{PENDING\_REVIEW}（待审核）。

%------------------------------------------------------------
\section{检测结果}

\subsection{数据有效性}

\begin{table}[h]
\centering
\begin{tabular}{lr}
\toprule
\textbf{项目} & \textbf{数量} \\
\midrule
总评论数         & 1000 \\
有效检测         & 813  \\
跳过（生成失败） & 38   \\
检测失败（网络）  & 147  \\
检测失败（其他） & 2    \\
\bottomrule
\end{tabular}
\caption{检测有效性统计}
\end{table}

共 187 条未获得有效结果，其中 147 条为网络连接错误（\texttt{Connection error}），
38 条为评论生成阶段失败，均已从指标统计中排除。

\subsection{混淆矩阵}

\begin{table}[h]
\centering
\begin{tabular}{lcc}
\toprule
 & \textbf{预测：剧透} & \textbf{预测：非剧透} \\
\midrule
\textbf{实际：剧透}   & TP = 264 & FN = 8 \\
\textbf{实际：非剧透} & FP = 2   & TN = 539 \\
\bottomrule
\end{tabular}
\caption{混淆矩阵（813 条有效样本）}
\end{table}

\subsection{性能指标}

\begin{table}[h]
\centering
\begin{tabular}{lr}
\toprule
\textbf{指标} & \textbf{数值} \\
\midrule
Accuracy（准确率）  & 98.8\% \\
Precision（精确率） & 99.2\% \\
Recall（召回率）    & 97.1\% \\
F1 Score            & 98.1\% \\
\bottomrule
\end{tabular}
\caption{剧透检测性能指标}
\end{table}

%------------------------------------------------------------
\section{结果分析}

\subsection{优势}

\begin{itemize}
  \item \textbf{极低误报率}（FP = 2，精确率 99.2\%）：系统几乎不会将正常评论误标为剧透，
        用户体验友好。
  \item \textbf{高召回率}（97.1\%）：仅有 8 条剧透评论漏报，
        剧透防护效果显著。
  \item \textbf{F1 = 98.1\%}：精确率与召回率均衡，整体性能优秀。
\end{itemize}

\subsection{不足与改进方向}

\begin{itemize}
  \item \textbf{网络稳定性}：147 条检测因网络错误失败（14.7\%），
        建议增加重试机制（exponential backoff）。
  \item \textbf{漏报分析}（FN = 8）：8 条剧透被判为非剧透，
        可能原因是该评论内容与后续段落语义距离较远，RAG 未召回相关上下文，
        LLM 仅凭常识判断时保守估计。可考虑扩大 top-k 或降低准入阈值。
  \item \textbf{测试集偏差}：评论由 Claude 生成，与真实用户评论风格存在分布差异，
        建议补充真实标注数据进行二次评估。
\end{itemize}

%------------------------------------------------------------
\section{结论}

在 813 条有效样本上，系统达到 \textbf{98.8\% 准确率}、\textbf{98.1\% F1}，
剧透检测效果优异，适合在阅读平台中作为自动审核的主要手段。
主要瓶颈在于网络连接稳定性，与模型判断能力无关。

\end{document}
